%! $n$th Roots \& Rational Exponents — Unit 5, Lesson 5.1 — Exit Ticket (Answer Key)
\documentclass[10pt]{article}
\usepackage{algebra2-article}
\usepackage{algebra2-key}   % pulls in -boxes; do NOT also load -boxes
\newcommand{\TallMath}[1]{$\displaystyle #1\rule[-1.4em]{0pt}{3.2em}$}

\begin{document}

\pageheader{Unit 5, Lesson 5.1}{Exit Ticket --- Answer Key}
\namedateperiod

\vspace{0.10in}

No notes. Assume all variables represent positive real numbers unless a question says otherwise.

\begin{enumerate}[leftmargin=1.8em, itemsep=10pt]
  \item \textbf{Translate and evaluate.}

        \vspace{4pt}
        (a) Write $\sqrt[5]{x^{3}}$ with a rational exponent: \ans{$x^{3/5}$}

        \vspace{4pt}
        (b) Write $y^{3/4}$ in radical form: \ans{$\sqrt[4]{y^{3}}$ \; (or $\left(\sqrt[4]{y}\right)^{3}$)}

        \vspace{4pt}
        (c) $27^{2/3} = $ \ans{$9$} \hspace{0.6in}
        (d) $16^{-1/2} = $ \ans{$\dfrac{1}{4}$}

        \vspace{4pt}
        (e) $x^{1/2}\cdot x^{1/6} = $ \ans{$x^{2/3}$}

  \item \textbf{Explain.} Using rational exponents, explain why $\sqrt[3]{-64}$ \emph{is} a real
        number but $\sqrt{-64}$ is \emph{not}. Your answer must mention the \textbf{denominator} of
        the exponent.
        \ansline{$\sqrt[3]{-64}=(-64)^{1/3}$. The denominator is $3$, which is \textbf{odd}, so it
        asks for a number whose \emph{cube} is $-64$ --- and cubing keeps the sign, so $(-4)^3=-64$
        works: the answer is $-4$.\\[3pt]
        $\sqrt{-64}=(-64)^{1/2}$. The denominator is $2$, which is \textbf{even}, so it asks for a
        number whose \emph{square} is $-64$. Every real number squared is positive or zero, so no
        such real number exists.}

  \item \textbf{SOL-style multiple choice.} Which expression is equivalent to $\sqrt[4]{x^{3}}$?
        \begin{enumerate}[label=\Alph*., leftmargin=1.9em, itemsep=1pt]
          \item $x^{4/3}$
          \item $x^{3/4}$ \quad \textcolor{keyred}{\textbf{$\leftarrow$ correct}}
          \item $x^{12}$
          \item $x^{-3/4}$
        \end{enumerate}
        \vspace{2pt}
        \begin{itemize}[leftmargin=1.4em, nosep]
          \item \textbf{B} is correct: the index $4$ is the \emph{denominator}, the power $3$ is the numerator.
          \item \textbf{A} is the signature error --- the fraction flipped.
          \item \textbf{C} multiplies the index by the power instead of dividing.
          \item \textbf{D} treats a root as a negative exponent; a root is a \emph{fractional} exponent, not a negative one.
        \end{itemize}
\end{enumerate}

\begin{teachernote}
Graded for completion --- mistakes happen, blanks don't. Sort the stack into four piles to decide
how Lesson 5.2 opens:
\begin{itemize}[leftmargin=1.3em, nosep]
  \item \textbf{Got it} --- 1(a)--(e) right and item 3 = B.
  \item \textbf{Flippers} --- item 3 = A, or 1(a) written as $x^{5/3}$. Re-run the
        $\left(a^{1/n}\right)^n=a$ derivation with them; memorizing ``index on the bottom'' without
        it does not survive a week.
  \item \textbf{Sign slips} --- 1(d) answered $-4$. One sentence fixes it: a negative exponent moves
        the expression, it does not change its sign.
  \item \textbf{Even/odd confusion} --- item 2 answered without reference to the denominator, or
        claiming both are real. This pile is the one that matters most: 5.2's product and quotient
        properties are stated \emph{only} for radicands that are legal, so a student who does not
        know when a root exists will simplify $\sqrt{-8}$ into existence.
\end{itemize}
Item 1(e) is the only one requiring fraction addition ($\frac12+\frac16=\frac23$); if that is where
the errors cluster, the problem is arithmetic, not radicals, and should be treated as such.
\end{teachernote}

\end{document}
